\documentclass[11pt,a4paper]{article}

% ── packages ─────────────────────────────────────────────────────────────────
\usepackage[margin=2.5cm]{geometry}
\usepackage{amsmath}
\usepackage{amssymb}
\usepackage{booktabs}
\usepackage{hyperref}
\usepackage{xcolor}
\usepackage{longtable}
\usepackage{array}
\usepackage{graphicx}
\usepackage{fancyhdr}
\usepackage{titlesec}
\usepackage{multirow}
\usepackage{enumitem}

% ── colours ──────────────────────────────────────────────────────────────────
\definecolor{apexblue}{RGB}{0,51,102}
\definecolor{apexgray}{RGB}{100,100,100}
\definecolor{apexred}{RGB}{180,30,30}

% ── hyperref setup ───────────────────────────────────────────────────────────
\hypersetup{
  colorlinks=true,
  linkcolor=apexblue,
  citecolor=apexblue,
  urlcolor=apexblue,
  pdftitle={MDD: PFE / CCR Exposure Engine -- APEX-MDL-0014 v1.1},
  pdfauthor={Apex Global Bank, Quantitative Research Division}
}

% ── header / footer ──────────────────────────────────────────────────────────
\pagestyle{fancy}
\fancyhf{}
\lhead{\textcolor{apexblue}{\small\textbf{APEX-MDL-0014} \textbar\ PFE / CCR Exposure Engine v1.1}}
\rhead{\textcolor{apexgray}{\small RESTRICTED --- Model Risk Sensitive}}
\lfoot{\textcolor{apexgray}{\small Apex Global Bank \textbar\ Quantitative Research Division}}
\rfoot{\textcolor{apexgray}{\small Page \thepage}}
\renewcommand{\headrulewidth}{0.4pt}
\renewcommand{\footrulewidth}{0.4pt}

% ── section style ────────────────────────────────────────────────────────────
\titleformat{\section}{\large\bfseries\color{apexblue}}{{\thesection.}}{0.5em}{}[\titlerule]
\titleformat{\subsection}{\normalsize\bfseries\color{apexblue}}{{\thesubsection}}{0.5em}{}
\titleformat{\subsubsection}{\normalsize\bfseries\color{apexgray}}{{\thesubsubsection}}{0.5em}{}

% ── misc ─────────────────────────────────────────────────────────────────────
\setlength{\parskip}{0.5em}
\setlength{\parindent}{0pt}

% ─────────────────────────────────────────────────────────────────────────────
\begin{document}
% ─────────────────────────────────────────────────────────────────────────────

% ── title block ──────────────────────────────────────────────────────────────
\begin{center}
  {\LARGE\bfseries\color{apexblue} Model Development Document}\\[0.4em]
  {\Large\bfseries PFE / CCR Exposure Engine}\\[0.3em]
  {\large Version 1.1 --- Revised per MRO Review Cycle 1}\\[0.6em]
  {\normalsize Apex Global Bank \textbar\ Quantitative Research Division}\\[0.2em]
  {\small\color{apexgray} SR 11-7 Section Reference: Model Development \textbar\ March 2026}
\end{center}

\vspace{1em}
\textcolor{apexred}{\textbf{Classification: RESTRICTED --- Model Risk Sensitive}}
\vspace{0.5em}

\textcolor{apexgray}{\small\textit{This revision addresses findings from the Model Risk Office
preliminary review of v1.0 (March 2026). Changed sections are marked \textbf{[REV]};
new sections are marked \textbf{[NEW]}.}}
\vspace{1em}

\hrule
\vspace{1em}

% ─────────────────────────────────────────────────────────────────────────────
\section{Document Metadata}
% ─────────────────────────────────────────────────────────────────────────────

\begin{table}[h!]
\centering
\begin{tabular}{@{}ll@{}}
\toprule
\textbf{Field} & \textbf{Value} \\
\midrule
Model Name           & PFE / CCR Exposure Engine \\
Model ID             & APEX-MDL-0014 \\
Version              & 1.1 \\
Status               & In Validation (Cycle 1 Response) \\
Risk Tier            & Tier 1 \\
Regulatory Framework & Basel III SA-CCR (BCBS 279); BCBS 261 (MPoR); IFRS 13; SR 11-7 \\
Date of Development  & March 2026 \\
Business Owner       & Head of Global Markets Trading (James Okafor) \\
Model Developer      & Quantitative Research --- XVA Team \\
MRO Review           & In Progress --- Dr.\ Rebecca Chen \\
CRO Approval         & Pending --- Dr.\ Priya Nair \\
\bottomrule
\end{tabular}
\caption{Document metadata}
\end{table}

\subsection{Version History}

\begin{table}[h!]
\centering
\begin{tabular}{@{}lll@{}}
\toprule
\textbf{Version} & \textbf{Date} & \textbf{Change Summary} \\
\midrule
0.1 & Jan 2026 & Initial netting framework; EPE/ENE profiles only \\
0.8 & Feb 2026 & Collateral offset with MPoR; SA-CCR add-on module \\
1.0 & Mar 2026 & Production candidate; submitted for validation \\
1.1 & Mar 2026 & MRO Cycle 1 response: real-world scaling defined; full SA-CCR \\
    &           & IR 3-bucket aggregation; VM mechanics; collateral simulation; \\
    &           & backtesting redesign; dispute MPoR trigger; stressed EPE; \\
    &           & factor space decomposition appendix; IMM gap analysis appendix \\
\bottomrule
\end{tabular}
\caption{Version history}
\end{table}

% ─────────────────────────────────────────────────────────────────────────────
\section{Executive Summary}
% ─────────────────────────────────────────────────────────────────────────────

\subsection{Model Purpose}

The PFE / CCR Exposure Engine is the foundational simulation model for all
counterparty credit risk (CCR) measurement at Apex Global Bank. It generates
the exposure profiles --- Expected Positive Exposure (EE/EPE), Expected Negative
Exposure (ENE), and Potential Future Exposure (PFE) --- consumed by three
downstream Tier~1 models:

\begin{itemize}[noitemsep]
  \item \textbf{CVA} (APEX-MDL-0010): uses EE profiles to compute the
        market value of counterparty default risk;
  \item \textbf{FVA} (APEX-MDL-0011): uses EE and ENE profiles to compute
        the funding cost of uncollateralised and partially-collateralised
        positions;
  \item \textbf{MVA} (APEX-MDL-0012): uses simulated future Initial Margin
        profiles derived from the same Monte Carlo paths.
\end{itemize}

In addition, the model produces:

\begin{itemize}[noitemsep]
  \item \textbf{PFE for credit limit management}: the 97.5th-percentile
        real-world exposure used by the Credit Risk function to set and
        monitor counterparty credit limits;
  \item \textbf{SA-CCR capital}: the Basel~III standardised-approach CCR
        Exposure-at-Default (EAD) used for regulatory capital calculation.
\end{itemize}

Because this model is upstream of multiple Tier~1 models and directly feeds
regulatory capital reporting, it is classified as \textbf{Tier~1}.

\subsection{Output Description}

\begin{table}[h!]
\centering
\begin{tabular}{@{}lll@{}}
\toprule
\textbf{Output} & \textbf{Consumer} & \textbf{Frequency} \\
\midrule
EE profile (time series, risk-neutral) & CVA, FVA, MVA models       & Daily \\
ENE profile (time series, risk-neutral) & FVA model                 & Daily \\
PFE at 97.5\% (time series, real-world) & Credit limit system       & Daily \\
Effective EPE (Basel 1-year integral)   & Regulatory capital (IMM)  & Daily \\
SA-CCR EAD by netting set               & Capital reporting, RWA    & Daily \\
Netting benefit report                  & Credit / legal            & Monthly \\
\bottomrule
\end{tabular}
\caption{Model outputs and consumers}
\end{table}

\subsection{Key Assumptions \textbf{[REV]}}

\begin{enumerate}[noitemsep]
  \item \textbf{Measure duality}: XVA outputs (EE, ENE for CVA/FVA/MVA) are
        computed under the risk-neutral measure $\mathbb{Q}$. PFE for credit
        limits is computed under the real-world measure $\mathbb{P}$, obtained
        by applying a deterministic real-world drift adjustment to the
        risk-neutral simulation. The conversion is fully specified in
        Section~\ref{sec:measure}.
  \item Market factor dynamics follow the parametric models specified in
        Section~\ref{sec:market_models} (Hull-White for rates, GBM for FX,
        Heston for equity, CIR for credit spreads).
  \item Netting sets are legally isolated: exposure netting is applied within
        but not across ISDA Master Agreements.
  \item Variation margin is called and received with a one-business-day
        settlement lag in normal conditions; collateral posted during the
        Margin Period of Risk is frozen at its last-received value
        (Section~\ref{sec:vm_mechanics}).
  \item SA-CCR calculation follows BCBS~279 (March 2014) as implemented in
        US rules effective April 2020, including the full three-maturity-bucket
        IR aggregation specified in BCBS~279 Annex~B
        (Section~\ref{sec:sa_ccr_full}).
\end{enumerate}

\subsection{Known Limitations}

\begin{enumerate}[noitemsep]
  \item The Gaussian copula used for joint factor simulation underestimates
        tail dependence in stress scenarios.
  \item The correlation matrix is static and recalibrated annually; intra-year
        correlation shifts are not captured.
  \item Gap risk for equity derivatives (overnight jumps) is not explicitly
        modelled.
  \item Approximately 8\% of counterparties require proxy CDS spread mapping,
        introducing basis risk.
  \item Path-dependent exotic derivatives use computational approximations that
        may introduce pricing errors of up to 5\%.
  \item \textbf{[NEW]} FX rates for long-dated positions ($>5$ years) are
        modelled with constant ATM volatility, omitting the FX volatility smile.
        See Section~\ref{sec:limits}.
  \item The intraday abbreviated run (500 paths) is insufficient for reliable
        PFE estimation at the 97.5th percentile; it is restricted to delta
        approximation of EE and not used for credit limit monitoring.
        See Section~\ref{sec:intraday}.
\end{enumerate}

% ─────────────────────────────────────────────────────────────────────────────
\section{Business Context and Purpose}
% ─────────────────────────────────────────────────────────────────────────────

\subsection{The Counterparty Credit Risk Problem}

OTC derivatives expose Apex to the risk that a counterparty defaults when the
derivative has positive mark-to-market value to Apex. Unlike a loan, where
exposure is the outstanding principal, a derivative's exposure is a random
variable that depends on the future evolution of market prices. The fundamental
challenge of CCR is that this exposure must be estimated over the entire life
of a contract --- potentially 30 years for a long-dated interest rate swap.

The 2008 financial crisis demonstrated the catastrophic consequences of
underestimating CCR. Bear Stearns and Lehman Brothers had large derivatives
portfolios with hundreds of counterparties. Exposure estimates that failed to
account for netting, or that used inadequate correlation assumptions, produced
dangerously low credit limits. When stress arrived, actual losses exceeded
modelled estimates by factors of 3--5x.

\subsection{Regulatory History}

CCR modelling has been shaped by successive waves of Basel regulation:

\begin{itemize}[noitemsep]
  \item \textbf{Basel II (2004)}: introduced the Internal Model Method (IMM)
        allowing banks with approved simulation models to use EPE for CCR
        capital, replacing the cruder Current Exposure Method (CEM).
  \item \textbf{Basel~III (2010)}: introduced the CVA capital charge requiring
        EPE-based CVA calculation; introduced the stressed EPE concept.
  \item \textbf{BCBS~279 / SA-CCR (2014)}: replaced CEM with the more
        risk-sensitive Standardised Approach to CCR, mandatory from 2020 for
        banks without IMM approval.
  \item \textbf{BCBS~261 (2013)}: established margin requirements for
        non-centrally-cleared derivatives, including MPoR specifications and
        dispute-history triggers.
  \item \textbf{FRTB-CVA (2017)}: reformed the CVA capital framework,
        requiring banks to compute CVA sensitivities using an approved
        simulation infrastructure.
\end{itemize}

Apex currently uses SA-CCR for regulatory capital and maintains this model as
a candidate for IMM approval. A gap analysis relative to IMM requirements is
provided in Appendix~\ref{app:imm_gap}.

\subsection{Model Users}

\begin{table}[h!]
\centering
\begin{tabular}{@{}lll@{}}
\toprule
\textbf{User} & \textbf{Function} & \textbf{Output Used} \\
\midrule
Credit Risk Officers   & Counterparty limit setting and monitoring & PFE at 97.5\% \\
XVA Desk               & CVA/FVA/MVA pricing and hedging           & EE, ENE \\
Regulatory Capital     & Basel III SA-CCR RWA calculation          & SA-CCR EAD \\
Finance / P\&L         & CVA P\&L via downstream CVA model         & EE \\
Internal Audit / MRO   & Model validation and ongoing monitoring   & All outputs \\
Regulators (OCC, Fed)  & CCAR stress testing; SR 11-7 review       & SA-CCR EAD, EE \\
\bottomrule
\end{tabular}
\caption{Model users and primary outputs consumed}
\end{table}

% ─────────────────────────────────────────────────────────────────────────────
\section{Theoretical Foundation}
\label{sec:theory}
% ─────────────────────────────────────────────────────────────────────────────

\subsection{Close-Out Netting}

Under an ISDA Master Agreement, all transactions between two parties within a
single netting set are aggregated upon default. The close-out netting value at
time $t$ is:

\begin{equation}
  V_{\text{net}}(t) = \sum_{k=1}^{K} V_k(t)
  \label{eq:netting}
\end{equation}

where $V_k(t)$ is the mark-to-market value of transaction $k$ at time $t$,
measured from Apex's perspective (positive = Apex is owed money).

The \emph{netting benefit} at time $t$ quantifies how much the netting agreement
reduces gross positive exposure:

\begin{equation}
  \text{NettingBenefit}(t)
    = \sum_{k=1}^{K} \max(V_k(t),\, 0)
    - \max\!\left(\sum_{k=1}^{K} V_k(t),\, 0\right)
  \label{eq:netting_benefit}
\end{equation}

A netting benefit of zero implies either all trades are in the same direction
(no offset) or the netting set is entirely out-of-the-money for Apex.

% ─────────────────────────────────────────────────────────────────────────────
\subsection{Variation Margin Mechanics and CSA Terms}
\label{sec:vm_mechanics}
% ─────────────────────────────────────────────────────────────────────────────

\subsubsection{Credit Support Annex Structure}

A Credit Support Annex (CSA), typically an ISDA 2016 Credit Support Deed
(CSD) for variation margin, governs the daily exchange of \emph{variation
margin} (VM) between counterparties. The key economic parameters are:

\begin{itemize}[noitemsep]
  \item \textbf{Threshold ($\Theta$)}: the amount of unsecured credit each
        party extends to the other before margin is called. Zero for
        fully-collateralised CSAs.
  \item \textbf{Minimum Transfer Amount (MTA, $\delta$)}: the minimum size of
        any single margin call, preventing operational burden from trivially
        small transfers (typically \$250K--\$500K).
  \item \textbf{Net Independent Collateral Amount (NICA)}: initial margin held
        (excluding VM), used in the RC floor of SA-CCR.
  \item \textbf{Eligible Collateral}: the schedule of securities and cash that
        may be posted as VM. Standard CSAs accept USD/EUR/GBP/JPY cash and
        government bonds subject to haircut.
  \item \textbf{Substitution Right}: the posting party may, on any business
        day, substitute existing posted collateral with alternative eligible
        collateral, provided the replacement value (after haircut) is
        equivalent.
  \item \textbf{Settlement Lag}: for securities collateral, $T+1$ or $T+2$
        standard settlement; for cash, same-day or $T+1$.
\end{itemize}

\subsubsection{Daily Variation Margin Call Process}

The VM call cycle operates as follows on each business day $t$:

\begin{enumerate}
  \item \textbf{Valuation}: The \emph{valuation agent} (per the CSA, typically
        the bank or determined by alphabetical convention) computes the net
        MtM of all trades in the netting set, $V_{\text{net}}(t)$.
  \item \textbf{Credit Support Balance}: Let $C(t)$ be the current collateral
        balance (positive = Apex holds collateral). The \emph{credit support
        amount} is:
        \begin{equation}
          \text{CSA}(t) = \max\!\bigl(V_{\text{net}}(t) - \Theta,\; \text{NICA}\bigr)
          \label{eq:csa_amount}
        \end{equation}
  \item \textbf{Delivery amount} (counterparty posts to Apex):
        \begin{equation}
          \text{DeliveryAmt}(t) =
          \begin{cases}
            \text{CSA}(t) - C(t) & \text{if } \text{CSA}(t) - C(t) \geq \delta \\
            0 & \text{otherwise}
          \end{cases}
          \label{eq:delivery}
        \end{equation}
  \item \textbf{Return amount} (Apex returns to counterparty):
        \begin{equation}
          \text{ReturnAmt}(t) =
          \begin{cases}
            C(t) - \text{CSA}(t) & \text{if } C(t) - \text{CSA}(t) \geq \delta \\
            0 & \text{otherwise}
          \end{cases}
          \label{eq:return}
        \end{equation}
  \item \textbf{Settlement}: Collateral moves with a $\Delta_{s}$-business-day
        settlement lag (typically 1 day for cash, 2 days for bonds). The
        updated balance is $C(t + \Delta_s) = C(t) + \text{DeliveryAmt}(t)
        - \text{ReturnAmt}(t)$.
\end{enumerate}

\subsubsection{Dispute Resolution}

If the counterparty disputes the valuation agent's MtM, the following process
governs:

\begin{enumerate}[noitemsep]
  \item The disputing party notifies the valuation agent by the applicable
        \emph{notification time} (typically noon local time).
  \item Both parties obtain independent mid-market quotes from at least two
        Reference Market Makers.
  \item If the parties agree within the \emph{Resolution Time} (typically 5
        business days), the agreed amount is exchanged.
  \item If no agreement is reached, the undisputed portion (the minimum of the
        two valuations net of any MTA) is transferred; the disputed remainder
        is held in escrow or the last agreed amount is maintained.
\end{enumerate}

Per BCBS~261 paragraph 49, if a netting set experiences \textbf{two or more
disputes within any rolling 24-month period} where the disputed amount exceeds
the greater of \$50M or 10\% of the threshold, the MPoR for that netting set
is doubled for a period of 24 months following the last dispute. This
\emph{dispute-history MPoR trigger} is explicitly tracked in the CSA database;
see Section~\ref{sec:mpor}.

\subsubsection{Collateral Haircuts}

Eligible securities are subject to supervisory haircuts to account for price
and FX risk during the MPoR. The collateral value after haircut is:

\begin{equation}
  C_{\text{net}}(t) = \sum_{i} C_i(t) \cdot (1 - H_i)(1 - H_{\text{FX},i})
  \label{eq:haircut}
\end{equation}

where $C_i(t)$ is the face value of security $i$, $H_i$ is the supervisory
security haircut (e.g.\ 2\% for 1--5 year government bonds), and
$H_{\text{FX},i} = 8\%$ if the security currency differs from the trade
currency. Haircut parameters follow BCBS~261 Table~1 and are input into the
CSA terms database.

% ─────────────────────────────────────────────────────────────────────────────
\subsection{Margin Period of Risk}
\label{sec:mpor}
% ─────────────────────────────────────────────────────────────────────────────

The Margin Period of Risk (MPoR) is the time from the last successful margin
call to the close-out of the portfolio following a counterparty default. During
MPoR, Apex cannot call additional margin; the last received collateral is
frozen. BCBS~261 specifies the following minimum MPoR values:

\begin{table}[h!]
\centering
\begin{tabular}{@{}ll@{}}
\toprule
\textbf{Condition} & \textbf{Minimum MPoR} \\
\midrule
Standard daily-margined netting set & 10 business days \\
Netting set with $\geq 5{,}000$ transactions & 20 business days \\
Netting set with illiquid or non-standard collateral & 20 business days \\
Re-hypothecated collateral & At least 20 business days \\
\textbf{[NEW]} Netting set with dispute-history trigger & $2\times$ applicable base MPoR \\
\bottomrule
\end{tabular}
\caption{MPoR specifications per BCBS~261}
\end{table}

The \textbf{dispute-history MPoR doubling} (BCBS~261 \S49--50) applies when a
netting set has experienced two or more material disputes within the preceding
24 months. The XVA Technology team maintains a dispute-history flag in the CSA
database, updated by the Collateral Management desk upon dispute resolution.
When the flag is set, the applicable MPoR for that netting set is doubled in
all downstream calculations including PFE, CVA, and SA-CCR.

The effective collateral at simulation time $t$, accounting for MPoR delay
$\tau$ (in simulation time units), is:

\begin{equation}
  C_{\text{eff}}(t) = C(t - \tau)
  \label{eq:coleff}
\end{equation}

The MPoR-adjusted collateralised exposure is therefore:

\begin{equation}
  V_{\text{col,MPoR}}(t)
    = \max\!\bigl(V_{\text{net}}(t) - C_{\text{net}}(t-\tau) - \Theta - \delta,\; 0\bigr)
  \label{eq:mpor_exposure}
\end{equation}

where $C_{\text{net}}$ is the haircut-adjusted balance from
Equation~\eqref{eq:haircut}.

% ─────────────────────────────────────────────────────────────────────────────
\subsection{Collateral Simulation on Monte Carlo Paths}
\label{sec:col_sim}
% ─────────────────────────────────────────────────────────────────────────────

\subsubsection{Rationale}

For a collateralised netting set, the collateral balance $C(t)$ is not a
fixed input --- it is itself a stochastic process driven by the same market
factors as $V_{\text{net}}(t)$. A margin call depends on today's MtM; the
collateral received depends on what MtM was $\tau$ days ago (MPoR delay).
Capturing this feedback loop correctly is essential for accurate collateralised
exposure simulation.

\subsubsection{Collateral Balance Recursion}

The collateral balance evolves on each simulation path $\omega$ according to
the following recursion. Let $t_i$ denote simulation time step $i$. Define
the \emph{target collateral} (the amount that would fully collateralise the
netting set given current MtM) as:

\begin{equation}
  C^*(t_i, \omega) = \max\!\bigl(V_{\text{net}}(t_i, \omega) - \Theta,\; \text{NICA}\bigr)
  \label{eq:col_target}
\end{equation}

The margin call triggered at time $t_i$ (delivery direction from counterparty
to Apex) is:

\begin{equation}
  \text{Call}(t_i, \omega) =
  \begin{cases}
    C^*(t_i, \omega) - C(t_{i-1}, \omega)
      & \text{if } C^*(t_i, \omega) - C(t_{i-1}, \omega) \geq \delta \\
    0 & \text{otherwise}
  \end{cases}
  \label{eq:margin_call}
\end{equation}

Settlements arrive at time $t_i + \Delta_s$, so the collateral balance updates
as:

\begin{equation}
  C(t_i, \omega) = C(t_{i-1}, \omega) + \text{Call}(t_{i-1} - \Delta_s, \omega)
  - \text{Return}(t_{i-1} - \Delta_s, \omega)
  \label{eq:col_balance}
\end{equation}

\subsubsection{MPoR Freezing}

In the PFE simulation, upon a hypothetical default at time $t$, collateral
calls issued during the MPoR are not honoured. The \emph{frozen collateral
balance} as of the last pre-MPoR call (time $t - \tau$) is used:

\begin{equation}
  C_{\text{frozen}}(t, \omega) = C(t - \tau, \omega)
  \label{eq:col_frozen}
\end{equation}

This implements the BCBS~261 convention that no new margin is received once
the close-out process begins. The MPoR period $[\,t-\tau, t\,]$ therefore
exposes Apex to the change in MtM over the MPoR uncollateralised:

\begin{equation}
  \text{MPoR exposure}(t, \omega)
    = \max\!\bigl(V_{\text{net}}(t,\omega) - C_{\text{frozen}}(t,\omega)
      - \Theta - \delta,\; 0\bigr)
  \label{eq:mpor_sim}
\end{equation}

\subsubsection{Treatment of Threshold and MTA}

Within the simulation, the threshold $\Theta$ and MTA $\delta$ are applied
per the CSA terms database:

\begin{itemize}[noitemsep]
  \item For \textbf{zero-threshold CSAs} ($\Theta = 0$, $\delta = 0$):
        equation~\eqref{eq:mpor_sim} reduces to
        $\max(V_{\text{net}} - C_{\text{frozen}}, 0)$, which reflects pure
        MPoR gap risk.
  \item For \textbf{threshold CSAs} ($\Theta > 0$): the threshold acts as
        a permanent unsecured credit line; the model correctly captures that
        Apex bears the first $\Theta$ of exposure regardless of
        collateral calls.
  \item For \textbf{non-zero MTA} ($\delta > 0$): MTA introduces a rounding
        band around the margin call threshold, slightly over-stating effective
        exposure. The simulation respects the MTA band exactly per
        equation~\eqref{eq:margin_call}.
\end{itemize}

\subsubsection{Re-Hypothecation Risk}

For counterparties where Apex re-hypothecates posted collateral (lending it
onward in the securities financing market), the MPoR applicable to both
directions is extended to at least 20 business days per BCBS~261. The
re-hypothecation flag is tracked in the CSA database.

% ─────────────────────────────────────────────────────────────────────────────
\subsection{Exposure Metrics}
\label{sec:metrics}
% ─────────────────────────────────────────────────────────────────────────────

Let $\tilde{V}(t, \omega)$ denote the MPoR-adjusted collateralised exposure
on path $\omega$ at time $t$ per Equation~\eqref{eq:mpor_sim}. All
exposure metrics derive from the distribution of $\tilde{V}(t)$ across paths.

\textbf{Terminology note}: this document distinguishes \emph{Expected
Exposure} (EE), a function of time, from \emph{Expected Positive Exposure}
(EPE), its time-average. Both are risk-neutral quantities used in XVA.
\emph{Potential Future Exposure} (PFE) is the high-confidence-level quantile
computed under the real-world measure and used for credit limit management.

\subsubsection{Expected Exposure}

\begin{align}
  \text{EE}(t)  &= \mathbb{E}^{\mathbb{Q}}\!\left[\max\!\bigl(\tilde{V}(t),\, 0\bigr)\right]
    \approx \frac{1}{N}\sum_{\omega=1}^{N} \max(\tilde{V}(t,\omega), 0)
  \label{eq:ee} \\[4pt]
  \text{ENE}(t) &= \mathbb{E}^{\mathbb{Q}}\!\left[\min\!\bigl(\tilde{V}(t),\, 0\bigr)\right]
    \approx \frac{1}{N}\sum_{\omega=1}^{N} \min(\tilde{V}(t,\omega), 0)
  \label{eq:ene}
\end{align}

\subsubsection{Expected Positive Exposure}

EPE is the time-average of EE over a one-year horizon, used as the primary
input to CVA:

\begin{equation}
  \text{EPE} = \frac{1}{T}\int_0^T \text{EE}(t)\,dt
    \approx \frac{1}{T}\sum_{i=1}^{M} \text{EE}(t_i)\,\Delta t_i
  \label{eq:epe}
\end{equation}

with $T = 1$ year and quadrature weights $\Delta t_i = t_i - t_{i-1}$.

\subsubsection{Effective EPE (Regulatory)}

For Basel IMM purposes, the Effective Expected Exposure is the non-decreasing
envelope of EE:

\begin{equation}
  \text{EffEE}(t) = \max\!\bigl(\text{EE}(t),\; \text{EffEE}(t - \Delta t)\bigr),
  \quad \text{EffEE}(0) = V_{\text{net}}(0)
  \label{eq:eff_ee}
\end{equation}

The Effective EPE (for regulatory capital under IMM) is:

\begin{equation}
  \text{EffEPE} = \frac{1}{1\text{Y}} \int_0^{1\text{Y}} \text{EffEE}(t)\, dt
  \label{eq:eff_epe}
\end{equation}

\subsubsection{Potential Future Exposure}

PFE at confidence level $\alpha$ at time horizon $t$ is the $\alpha$-quantile
of the real-world exposure distribution (see Section~\ref{sec:measure} for
real-world conversion):

\begin{equation}
  \text{PFE}_\alpha(t) = Q_\alpha^{\mathbb{P}}\!\left[\tilde{V}(t)\right]
  \label{eq:pfe}
\end{equation}

Apex uses $\alpha = 97.5\%$. The distinction from
$Q_\alpha^{\mathbb{Q}}$ is material for long-dated positions and is
quantified in Section~\ref{sec:measure}.

% ─────────────────────────────────────────────────────────────────────────────
\subsection{Real-World vs Risk-Neutral Measure}
\label{sec:measure}
% ─────────────────────────────────────────────────────────────────────────────

\subsubsection{Motivation}

XVA quantities (CVA, FVA, MVA) are pricing adjustments consistent with
market-implied no-arbitrage dynamics and are computed under the risk-neutral
measure $\mathbb{Q}$. Credit limit management (PFE) is a risk management
exercise concerned with the actual distribution of future exposures; it should
use the real-world measure $\mathbb{P}$.

The difference is material: under $\mathbb{Q}$, interest rate drifts are
determined by the no-arbitrage condition (drift $= r_d - r_f$ for FX); under
$\mathbb{P}$, drifts reflect historical term premia and mean-reversion toward
long-run equilibrium levels (e.g.\ the 10-year USD SOFR rate averages 4.5\%
over the historical estimation window, whereas the current forward curve
implies a lower level for the next 5 years).

\subsubsection{Real-World Drift Adjustments}

The model uses a \emph{deterministic drift overlay} approach: paths are
generated under $\mathbb{Q}$ as the base case, and a time-dependent drift
adjustment $\mu^{\mathbb{P}}(t) - \mu^{\mathbb{Q}}(t)$ is applied to
obtain $\mathbb{P}$-measure paths. For the Hull-White short-rate model:

\begin{equation}
  r^{\mathbb{P}}(t) = r^{\mathbb{Q}}(t) + \phi(t)
  \label{eq:rw_adjust}
\end{equation}

where $\phi(t)$ is the \emph{risk premium function}, estimated from the
mean difference between historical realised short rates and the rates implied
by the $\mathbb{Q}$ forward curve over the calibration period 2015--2025:

\begin{equation}
  \hat{\phi}(t) = \frac{1}{n}\sum_{s=1}^{n} \bigl(r^{\text{hist}}(s+t) - f(s, t)\bigr)
  \label{eq:phi_est}
\end{equation}

where $f(s, t)$ is the $t$-year forward rate prevailing at historical date $s$.

For FX rates, the $\mathbb{P}$-measure drift incorporates the historical
excess return of foreign currency investment over USD:

\begin{equation}
  \frac{dS^{\mathbb{P}}(t)}{S(t)}
    = \bigl(r_d(t) - r_f(t) + \lambda_{\text{FX}}\bigr)\,dt
    + \sigma_{\text{FX}}\,dW_{\text{FX}}^{\mathbb{P}}(t)
  \label{eq:fx_rw}
\end{equation}

where $\lambda_{\text{FX}}$ is the FX risk premium (estimated at
$+0.8\%$ annualised for EUR/USD over the 2000--2025 window).

\subsubsection{PFE Real-World Scaling Factor}

The PFE scaling factor $\xi_{\text{RW}}$ is defined as the ratio of the
97.5th-percentile exposure under $\mathbb{P}$ to the 97.5th-percentile under
$\mathbb{Q}$ at the 1-year horizon, computed on the current portfolio:

\begin{equation}
  \xi_{\text{RW}} = \frac{Q_{97.5\%}^{\mathbb{P}}\!\left[\tilde{V}(1\text{Y})\right]}
                         {Q_{97.5\%}^{\mathbb{Q}}\!\left[\tilde{V}(1\text{Y})\right]}
  \label{eq:rw_scale}
\end{equation}

Current estimated value: $\xi_{\text{RW}} = 1.12$ (i.e.\ real-world PFE is
approximately 12\% higher than risk-neutral PFE at the 1-year horizon). The
scaling factor is re-estimated quarterly using a 200-path real-world
simulation and applied to the full 10,000-path risk-neutral PFE profile
to obtain the reported credit-limit PFE.

\textbf{Limitation}: the scaling factor is applied as a time-independent
scalar, whereas the true ratio $Q_{97.5\%}^{\mathbb{P}} / Q_{97.5\%}^{\mathbb{Q}}$
is likely horizon-dependent. Full real-world simulation with 10,000 paths is
planned for v2.0 (see Appendix~\ref{app:imm_gap}).

% ─────────────────────────────────────────────────────────────────────────────
\subsection{Market Factor Models}
\label{sec:market_models}
% ─────────────────────────────────────────────────────────────────────────────

\subsubsection{Interest Rates: Two-Factor Hull-White}

Domestic and foreign short rates are modelled using the two-factor Hull-White
(HW2) model, which fits the initial term structure exactly while allowing
mean-reverting stochastic evolution:

\begin{align}
  r(t) &= x_1(t) + x_2(t) + \varphi(t) \label{eq:hw2_r}\\
  dx_1(t) &= -a_1\, x_1(t)\, dt + \sigma_1\, dW_1(t) \label{eq:hw2_x1}\\
  dx_2(t) &= -a_2\, x_2(t)\, dt + \sigma_2\, dW_2(t) \label{eq:hw2_x2}
\end{align}

where $\varphi(t)$ is the deterministic fit to the initial forward curve,
$a_1, a_2$ are mean-reversion speeds, $\sigma_1, \sigma_2$ are volatilities,
and $\text{Corr}(dW_1, dW_2) = \rho_{12}$.

\textbf{Factor count}: each HW2 curve contributes \textbf{two} stochastic
factors $(x_1, x_2)$ to the joint simulation. With five rate curves (SOFR,
ESTR, SONIA, TONAR, HIBOR), the IR factor block accounts for
$5 \times 2 = 10$ factors. See Appendix~\ref{app:factors} for the complete
factor decomposition.

Current calibration (March 2026):
$a_1 = 0.03$, $a_2 = 0.35$, $\sigma_1 = 0.007$, $\sigma_2 = 0.014$,
$\rho_{12} = -0.72$ (applied uniformly across curves; curve-specific calibrations
are tracked in the calibration ledger).

\subsubsection{FX Rates: Geometric Brownian Motion}

FX spot rates $S_{\text{CCY/USD}}(t)$ are modelled as log-normal:

\begin{equation}
  \frac{dS(t)}{S(t)} = \bigl(r_d(t) - r_f(t)\bigr)\, dt + \sigma_{\text{FX}}\, dW_{\text{FX}}(t)
  \label{eq:fx_gbm}
\end{equation}

where $r_d$ and $r_f$ are domestic and foreign short rates from the HW2 models,
and $\sigma_{\text{FX}}$ is calibrated to ATM options with maturity matching
the simulation horizon.

\textbf{Limitation}: the constant-volatility GBM does not capture the FX
volatility smile, which is material for strikes far from spot. For long-dated
FX derivatives ($>5$ years), this can understate PFE at the 97.5th percentile
by an estimated 8--15\% depending on strike. This is listed as a known
limitation (see Section~\ref{sec:limits}) and addressed by the compensating
vol-of-vol add-on for affected trades.

\subsubsection{Equity: Heston Stochastic Volatility}

Equity prices are modelled under the Heston (1993) framework to capture the
volatility smile:

\begin{align}
  \frac{dS(t)}{S(t)} &= \mu\, dt + \sqrt{v(t)}\, dW_S(t) \label{eq:heston_s}\\
  dv(t) &= \kappa\bigl(\bar{v} - v(t)\bigr)\, dt + \xi\sqrt{v(t)}\, dW_v(t) \label{eq:heston_v}
\end{align}

with $\text{Corr}(dW_S, dW_v) = \rho_{\text{SV}}$. Parameters are calibrated
to the equity volatility surface. \textbf{The Heston model contributes two
stochastic factors per equity index} (log-price and variance), yielding
$12 \times 2 = 24$ equity factors in the joint simulation.

\subsubsection{Credit Spreads: Cox-Ingersoll-Ross}

Credit spreads follow a CIR process, ensuring non-negativity:

\begin{equation}
  d\lambda(t) = \kappa_\lambda\!\bigl(\bar{\lambda} - \lambda(t)\bigr)\, dt
               + \sigma_\lambda\sqrt{\lambda(t)}\, dW_\lambda(t)
  \label{eq:cir}
\end{equation}

The CIR contributes one factor per credit spread curve.

% ─────────────────────────────────────────────────────────────────────────────
\subsection{Joint Simulation via Cholesky Decomposition}
\label{sec:cholesky}
% ─────────────────────────────────────────────────────────────────────────────

All market factors are simulated jointly using a correlation matrix
$\mathbf{P}$, Cholesky-decomposed into $\mathbf{L}$ such that
$\mathbf{L}\mathbf{L}^\top = \mathbf{P}$. Standard normal variates
$\mathbf{z} \sim \mathcal{N}(\mathbf{0}, \mathbf{I})$ are transformed into
correlated variates $\mathbf{w} = \mathbf{L}\mathbf{z}$.

\textbf{[REV] Corrected factor decomposition}: the joint factor space has
$\mathbf{59}$ factors, not 47 as stated in v1.0. The revision arises from
correctly counting two factors per HW2 curve and two per Heston index.
The complete breakdown is in Appendix~\ref{app:factors}; the Cholesky
decomposition is therefore $59 \times 59$. The production code was already
implemented with 59 factors; this was a documentation error in v1.0.

% ─────────────────────────────────────────────────────────────────────────────
\subsection{SA-CCR: Full Standardised Approach Methodology}
\label{sec:sa_ccr_full}
% ─────────────────────────────────────────────────────────────────────────────

\subsubsection{EAD Formula}

For regulatory capital purposes under Basel~III (BCBS~279), the
Exposure-at-Default (EAD) for a netting set is:

\begin{equation}
  \text{EAD} = \alpha \cdot \bigl(\text{RC} + \text{PFE}_{\text{addon}}\bigr)
  \label{eq:sa_ccr_ead}
\end{equation}

where $\alpha = 1.4$ is the supervisory scalar.

\subsubsection{Replacement Cost}

For \textbf{margined} netting sets:

\begin{equation}
  \text{RC} = \max\!\bigl(V - C_{\text{net}},\; \text{TH} + \text{MTA} - \text{NICA},\; 0\bigr)
  \label{eq:rc_margined}
\end{equation}

For \textbf{unmargined} netting sets:

\begin{equation}
  \text{RC} = \max\!\bigl(V - C_{\text{net}},\; 0\bigr)
  \label{eq:rc_unmargined}
\end{equation}

where $V$ is current net MtM, $C_{\text{net}}$ is collateral held net of
haircuts per Equation~\eqref{eq:haircut}, TH is the CSA threshold, MTA is the
minimum transfer amount, and NICA is the net independent collateral amount.

\subsubsection{PFE Add-On: Aggregation Structure}

The aggregate PFE add-on sums independently across five asset classes:

\begin{equation}
  \text{PFE}_{\text{addon}} = \text{AddOn}^{\text{IR}}
    + \text{AddOn}^{\text{FX}}
    + \text{AddOn}^{\text{credit}}
    + \text{AddOn}^{\text{EQ}}
    + \text{AddOn}^{\text{com}}
  \label{eq:pfe_addon}
\end{equation}

No netting applies across asset classes.

\subsubsection{Interest Rate Add-On: Full BCBS 279 Annex B}
\label{sec:ir_addon}

For each currency \textit{ccy}, trades are assigned to one of three
\emph{maturity buckets} based on residual maturity $M_j$:

\begin{table}[h!]
\centering
\begin{tabular}{@{}cll@{}}
\toprule
\textbf{Bucket} & \textbf{Residual Maturity} & \textbf{Bucket Correlations} \\
\midrule
1 & $M_j \leq 1$ year & \\
2 & $1 < M_j \leq 5$ years & $\rho_{12} = \rho_{23} = 0.70$ \\
3 & $M_j > 5$ years & $\rho_{13} = 0.30$ \\
\bottomrule
\end{tabular}
\caption{SA-CCR IR maturity buckets and correlations (BCBS 279 Annex B)}
\end{table}

\paragraph{Adjusted Notional.}
For an interest rate derivative $j$, the adjusted notional captures the
dollar-value sensitivity of the trade to rate moves:

\begin{equation}
  d_j = N_j \cdot \text{SD}_j,
  \quad
  \text{SD}_j = \frac{e^{-0.05 S_j} - e^{-0.05 E_j}}{0.05}
  \label{eq:adj_notional}
\end{equation}

where $N_j$ is the trade notional, and $S_j$ and $E_j$ are the start and end
dates of the underlying rate (in years from the calculation date). For vanilla
swaps, $S_j = 0$ and $E_j$ is the swap maturity.

\paragraph{Maturity Factor.}
For unmargined netting sets:

\begin{equation}
  \text{MF}_j^{\text{unmargined}} = \sqrt{\frac{\min(M_j,\; 1\text{Y})}{1\text{Y}}}
  \label{eq:mf_unmargined}
\end{equation}

For margined netting sets, the maturity factor reflects the MPoR $\tau$
(converted to years):

\begin{equation}
  \text{MF}_j^{\text{margined}} = 1.5 \cdot \sqrt{\frac{\tau}{1\text{Y}}}
  \label{eq:mf_margined}
\end{equation}

where the factor 1.5 is the BCBS~279 supervisory scaling applied to all
margined netting sets.

\paragraph{Effective Notional per Bucket.}
The bucket effective notional, combining supervisory delta, adjusted notional,
maturity factor, and supervisory IR volatility $\sigma^{\text{IR}} = 0.50\%$:

\begin{equation}
  D_b^{\text{ccy}} = \sum_{j \in \text{bucket }b}
    \delta_j \cdot d_j \cdot \text{MF}_j \cdot \sigma^{\text{IR}}
  \label{eq:bucket_notional}
\end{equation}

where $\delta_j \in \{+1,\, -1\}$ is the supervisory delta (sign of DV01).
For options, $\delta_j$ is the Black-Scholes delta adjusted by sign, bounded
to $[-1, +1]$.

\paragraph{Currency Add-On (Three-Bucket Aggregation).}
The add-on for a single currency hedging set aggregates across buckets with
the correlation structure from BCBS~279 Annex~B \S4:

\begin{equation}
  \text{AddOn}^{\text{IR,ccy}} = \sqrt{
    \bigl(D_1^{\text{ccy}}\bigr)^2
    + \bigl(D_2^{\text{ccy}}\bigr)^2
    + \bigl(D_3^{\text{ccy}}\bigr)^2
    + 2\rho_{12}\, D_1^{\text{ccy}} D_2^{\text{ccy}}
    + 2\rho_{23}\, D_2^{\text{ccy}} D_3^{\text{ccy}}
    + 2\rho_{13}\, D_1^{\text{ccy}} D_3^{\text{ccy}}
  }
  \label{eq:ir_addon_full}
\end{equation}

with $\rho_{12} = \rho_{23} = 0.70$ and $\rho_{13} = 0.30$.

\paragraph{Full IR Add-On.}
Across all currencies (no netting across currencies):

\begin{equation}
  \text{AddOn}^{\text{IR}} = \sum_{\text{ccy}} \text{AddOn}^{\text{IR,ccy}}
  \label{eq:ir_total}
\end{equation}

\subsubsection{FX Add-On}

For each FX pair (currency pair CCY/USD treated as a single hedging set):

\begin{equation}
  \text{AddOn}^{\text{FX,pair}} = \left|\sum_j \delta_j \cdot d_j \cdot \text{MF}_j
    \cdot \sigma^{\text{FX}}\right|
  \label{eq:fx_addon}
\end{equation}

where $\sigma^{\text{FX}} = 4\%$ (supervisory volatility for FX).
No netting across pairs:
$\text{AddOn}^{\text{FX}} = \sum_{\text{pairs}} \text{AddOn}^{\text{FX,pair}}$.

\subsubsection{Credit Add-On}

For each reference entity, trades are grouped into a single-name hedging set.
The single-name add-on uses the same effective-notional approach with
$\sigma^{\text{credit}} \in \{0.38\%\text{ (IG)},\; 1.06\%\text{ (HY)}\}$
and bucket correlations of 0.35 within IG or HY and 0.0 across.

Index trades are treated as separate hedging sets. The index add-on uses
$\sigma^{\text{index}} \in \{0.20\%\text{ (IG index)},\; 0.60\%\text{ (HY index)}\}$.

\subsubsection{Equity Add-On}

Single-name equity derivatives use $\sigma^{\text{EQ}} = 32\%$. Index
derivatives use $\sigma^{\text{EQ,index}} = 20\%$. The partial netting
formula within a single-issuer hedging set applies correlation $\rho = 0.50$
across single names.

\subsubsection{Commodity Add-On}

Commodity derivatives are grouped by sub-class (energy, metals, agriculture,
other). Within a sub-class, the add-on formula is:

\begin{equation}
  \text{AddOn}^{\text{com,sub}} = \sqrt{
    \left(\sum_j \delta_j d_j \text{MF}_j \sigma^{\text{com}}\right)^2
    + \bigl(1-\rho_{\text{com}}^2\bigr)
      \sum_j \bigl(\delta_j d_j \text{MF}_j \sigma^{\text{com}}\bigr)^2
  }
  \label{eq:com_addon}
\end{equation}

with $\rho_{\text{com}} = 0.40$ for within-sub-class correlation and
$\sigma^{\text{com}} \in \{40\%\text{ (energy)},\; 18\%\text{ (metals)},
\; 18\%\text{ (agriculture)},\; 18\%\text{ (other)}\}$.

% ─────────────────────────────────────────────────────────────────────────────
\section{Data Requirements and Governance}
% ─────────────────────────────────────────────────────────────────────────────

\subsection{Input Data Sources}

\begin{longtable}{@{}lll@{}}
\toprule
\textbf{Data Item} & \textbf{Source} & \textbf{Frequency} \\
\midrule
\endhead
\bottomrule
\endfoot
OIS/RFR yield curves (SOFR, ESTR, SONIA, TONAR, HIBOR) & Bloomberg / Internal & Daily \\
FX spot rates (10 pairs vs.\ USD) & Bloomberg & Real-time (EOD snap) \\
Equity index levels (12 indices) & Bloomberg & Daily \\
Credit spread curves (IG, HY, fins, sov) & Markit / Bloomberg & Daily \\
Swaption vol surfaces (all rate curves) & Bloomberg & Daily \\
FX ATM vol (all pairs) & Bloomberg & Daily \\
Equity vol surfaces (all indices) & Bloomberg & Daily \\
Credit spread vol (CDS options where available) & Markit & Weekly \\
ISDA Master Agreement database & Legal / ISDA operations & On amendment \\
CSA terms (threshold, MTA, eligible collateral, haircuts) & Legal / ISDA operations & On amendment \\
Dispute-history flags (BCBS 261 MPoR trigger) & Collateral Management & On update \\
Trade population (all OTC derivatives) & Front-office systems & Real-time \\
Collateral balances (VM posted/received) & Collateral Management & Daily \\
Historical factor data (5-year lookback) & Bloomberg / Proprietary & Annual \\
Real-world risk premium estimates $\hat{\phi}(t)$ & Internal QR & Quarterly \\
\end{longtable}

\subsection{Data Quality Controls}

\begin{enumerate}[noitemsep]
  \item \textbf{Curve completeness}: All OIS curves must have at least 10 tenor
        points covering 1W to 30Y. Missing tenors use linear log-discount
        interpolation with a staleness flag if the bracketing points are
        $>2$ business days old.
  \item \textbf{Vol surface consistency}: Swaption vol surfaces are validated
        for absence of calendar spread arbitrage and butterfly spread arbitrage
        before use. Surfaces failing validation revert to prior-day with an
        alert.
  \item \textbf{Trade data completeness}: Trade population must reconcile to
        front-office systems within 0.5\% of notional before the simulation
        run proceeds. Material breaks ($>1\%$) halt the run and escalate to
        Operations.
  \item \textbf{CSA database}: The CSA terms database must cover $\geq 99\%$
        of live netting sets by notional. Missing CSA terms default to
        ``uncollateralised'' (conservative).
  \item \textbf{[NEW] Dispute-history flags}: The collateral management system
        must report dispute-history flags for all netting sets. The flag is
        set by the Collateral Management desk within 1 business day of dispute
        resolution. The MPoR override table is reconciled to the CSA database
        on a weekly basis.
\end{enumerate}

\subsection{BCBS 239 Compliance}

As a Tier~1 model feeding regulatory capital, the PFE/CCR engine is subject to
BCBS~239 data lineage requirements. Every input data item is tagged with
provenance metadata (source system, timestamp, transformation applied).
The full data lineage from raw market data to SA-CCR EAD must be
reconstructable within 4 hours of a regulatory request.

% ─────────────────────────────────────────────────────────────────────────────
\section{Methodology and Implementation}
% ─────────────────────────────────────────────────────────────────────────────

\subsection{Monte Carlo Engine}
\label{sec:mc}

The simulation engine generates $N = 10{,}000$ risk-neutral paths over
$M = 200$ time steps using the following time grid:

\begin{itemize}[noitemsep]
  \item $t_1, \ldots, t_{52}$: weekly steps to 1 year;
  \item $t_{53}, \ldots, t_{112}$: monthly steps from 1 to 5 years;
  \item $t_{113}, \ldots, t_{200}$: quarterly steps from 5 to 30 years.
\end{itemize}

\textbf{[NEW] Time grid adequacy for short-dated instruments}: the weekly
grid from 0 to 1 year provides 52 time steps at 5-business-day resolution.
For instruments with remaining maturity $<3$ months at a simulation date,
the weekly grid may under-sample the exposure profile near maturity. The
SA-CCR maturity factor $\text{MF}_j = \sqrt{M_j / 1\text{Y}}$ implicitly
accounts for this by weighting short-dated positions lower. However, for
PFE estimation of short-dated options (e.g.\ 1-month FX options), the
discrete 5-day steps may produce a stepped exposure profile. A daily
sub-grid for instruments with residual maturity $< 3$ months is planned
for v2.0.

GPU acceleration (CUDA via cuBLAS) is used for the $59 \times 59$ Cholesky
decomposition and for batched path generation. The full portfolio run
completes in under 4 hours on the production GPU cluster.

\subsection{Intraday Abbreviated Run and Path Count Adequacy}
\label{sec:intraday}

An intraday ``Greeks'' run uses 500 paths over a 1-year horizon only.
\textbf{[NEW] Statistical adequacy assessment}: at $N = 500$ paths, the
97.5th percentile is estimated from approximately $500 \times 0.025 = 12.5$
observations above the threshold. The standard error of the quantile
estimate is approximately:

\begin{equation}
  \text{SE}\!\left[\hat{Q}_{97.5\%}\right]
    \approx \frac{\sqrt{p(1-p)}}{n \cdot f(\hat{Q}_{97.5\%})} \cdot \sigma
    \approx \frac{\sqrt{0.025 \times 0.975}}{500 \cdot f(\hat{Q}_{97.5\%})}
  \label{eq:se_pfe}
\end{equation}

Empirical testing shows the 90\% confidence interval around the 500-path
PFE$_{97.5\%}$ spans $\pm 35\%$ of the full-run estimate. As a result:

\begin{itemize}[noitemsep]
  \item The intraday 500-path run is \textbf{restricted} to delta-approximation
        of EE only (using analytical Greeks) and is \textbf{not used} for
        PFE estimation or credit limit monitoring.
  \item Credit limit monitoring relies solely on the prior-day full 10,000-path
        PFE run. Intraday credit limit breaches trigger the Greek-scaled prior-day
        PFE as the interim estimate, clearly labelled as approximate.
\end{itemize}

\subsection{Collateral Simulation Algorithm}
\label{sec:col_algo}

The following algorithm implements Section~\ref{sec:col_sim} on each Monte
Carlo path $\omega$:

\begin{enumerate}
  \item \textbf{Initialise}: set $C(t_0, \omega) = C_{\text{market}}$ (current
        actual collateral balance from the Collateral Management system).
  \item \textbf{For each time step} $t_i$, $i = 1, \ldots, M$:
    \begin{enumerate}
      \item Compute $V_{\text{net}}(t_i, \omega)$ by pricing all trades under
            simulated market factors at $t_i$.
      \item Compute the CSA target $C^*(t_i, \omega)$ per
            Equation~\eqref{eq:col_target}.
      \item Compute $\text{Call}(t_i, \omega)$ per Equation~\eqref{eq:margin_call}.
      \item Update collateral balance: $C(t_i, \omega) = C(t_{i-1}, \omega)
            + \text{Call}(t_{i-\Delta_s}, \omega) - \text{Return}(t_{i-\Delta_s}, \omega)$
            (incorporating settlement lag $\Delta_s$, typically $i - 1$ for a
            1-step lag on the weekly grid).
      \item Apply MPoR freeze: $C_{\text{frozen}}(t_i, \omega) = C(t_{i - \tau_{\text{grid}}}, \omega)$
            where $\tau_{\text{grid}}$ is the number of time steps corresponding
            to the applicable MPoR (rounded up to the next step boundary).
      \item Compute $\tilde{V}(t_i, \omega)$ per Equation~\eqref{eq:mpor_sim}.
    \end{enumerate}
  \item \textbf{Aggregate}: after all paths, compute EE, ENE, PFE per
        Section~\ref{sec:metrics}.
\end{enumerate}

\textbf{Boundary condition}: at $t_0$, the actual collateral balance is used
rather than a simulated value. For multi-path consistency, all paths start from
the same observed $C(t_0)$.

\textbf{Re-hypothecation paths}: for netting sets with a re-hypothecation flag,
$\tau_{\text{grid}}$ is set to the 20-business-day MPoR minimum regardless of
transaction count.

\subsection{SA-CCR Calculation}

The SA-CCR module runs as a post-processing step on the current-date trade
population (no simulation paths required). For each netting set:

\begin{enumerate}[noitemsep]
  \item Classify each trade into asset class (IR, FX, credit, EQ, commodity).
  \item Compute supervisory delta $\delta_j$ for each trade.
  \item Compute adjusted notional $d_j$ per Equation~\eqref{eq:adj_notional}
        (tenor-weighted supervisory duration for IR; notional for FX/EQ/credit).
  \item Compute maturity factor $\text{MF}_j$ per
        Equations~\eqref{eq:mf_unmargined} or \eqref{eq:mf_margined}.
  \item For IR: assign to maturity bucket (1/2/3) and aggregate per
        Equation~\eqref{eq:ir_addon_full}.
  \item For all other asset classes: aggregate per respective formulas.
  \item Sum to $\text{PFE}_{\text{addon}}$ per Equation~\eqref{eq:pfe_addon}.
  \item Compute RC per Equation~\eqref{eq:rc_margined} or \eqref{eq:rc_unmargined}.
  \item Compute EAD per Equation~\eqref{eq:sa_ccr_ead}.
\end{enumerate}

\subsection{Stressed EPE Methodology}
\label{sec:stressed_epe}

\textbf{[NEW]} Stressed EPE is computed quarterly using two methodologies:

\paragraph{Scenario-Based Stressed EPE.}
Five named macro scenarios (see Table in Section~\ref{sec:stress_scenarios})
are applied as instantaneous shocks to the initial market state. The full
Monte Carlo simulation then runs from the shocked initial conditions.
Stressed EE$(t)$ and stressed EPE are computed identically to the base case,
with the shocked state as $t_0$.

\paragraph{Stressed-Period Calibrated EPE.}
Following Basel III \S\S\,277--278, stressed EPE uses model parameters
calibrated to a ``significantly stressed'' historical period. The calibration
period identified as most stressed for the current portfolio is:
\textbf{1 October 2008 -- 30 September 2009} (Lehman/GFC). Stressed calibration
replaces the base HW2/Heston/CIR parameters with those estimated from this
period, holding the current MtM and trade population fixed.

The ratio of stressed EPE to base EPE (the \emph{stress multiplier}) is:

\begin{equation}
  m_{\text{stress}} = \frac{\text{EPE}_{\text{stressed}}}{\text{EPE}_{\text{base}}}
  \label{eq:stress_mult}
\end{equation}

Current estimated stress multiplier: $m_{\text{stress}} = 1.42$ (March 2026
portfolio). The stressed EPE is reported to the Market Risk Officer quarterly
and disclosed in the Board Risk Committee pack.

\subsection{Shared Infrastructure with CVA, FVA, MVA}

CVA, FVA, and MVA all consume outputs from the same Monte Carlo run, ensuring
consistency under IFRS~13. The dependency chain is:

\begin{center}
  \textbf{PFE/CCR Engine} $\;\longrightarrow\;$ \textbf{CVA} (APEX-MDL-0010)\\[2pt]
  \textbf{PFE/CCR Engine} $\;\longrightarrow\;$ \textbf{FVA} (APEX-MDL-0011)\\[2pt]
  \textbf{PFE/CCR Engine} $\;\longrightarrow\;$ \textbf{MVA} (APEX-MDL-0012)
\end{center}

\textbf{[REV] Time step consistency}: CVA (APEX-MDL-0010) processes the same
200-step EE profile generated by this engine. The v1.0 CVA documentation
referred to ``30 time steps'' in its methodology section; this is an error in
the CVA MDD, which will be corrected in a forthcoming CVA v1.1 revision. The
authoritative time grid is the 200-step grid defined in this document.

A failure of the PFE/CCR engine propagates to all three downstream models.
The fallback procedure (prior-day EPE/ENE scaled by Greeks) is documented
in Section~\ref{sec:controls}.

% ─────────────────────────────────────────────────────────────────────────────
\section{Model Testing and Backtesting}
% ─────────────────────────────────────────────────────────────────────────────

\subsection{P\&L Explain for EE}

Daily predicted EE change is compared to actual EE change:

\begin{equation}
  \Delta\text{EE}_{\text{predicted}}
    = \frac{\partial \text{EE}}{\partial r} \Delta r
    + \frac{\partial \text{EE}}{\partial S_{\text{FX}}} \Delta S_{\text{FX}}
    + \frac{\partial \text{EE}}{\partial \lambda} \Delta\lambda
    + \text{theta}
\end{equation}

P\&L explain ratio $= \Delta\text{EE}_{\text{predicted}} / \Delta\text{EE}_{\text{actual}}$
must exceed 85\%. Deviations $>15\%$ trigger a model investigation within
2 business days.

\subsection{SA-CCR Benchmark}

The SA-CCR EAD is compared quarterly to the 97.5th-percentile internal PFE
multiplied by $\alpha = 1.4$:

\begin{equation}
  \text{SA-CCR ratio}
    = \frac{\text{SA-CCR EAD}}{1.4 \times \text{PFE}_{97.5\%}(1\text{Y})}
\end{equation}

Acceptable range: $[0.9,\; 1.5]$. Values outside this range indicate either
over-conservatism in SA-CCR or inadequacy of the internal model, and trigger
a 30-day review.

\subsection{Historical Backtesting}
\label{sec:backtest}

\textbf{[REV] Revised backtesting design.} The v1.0 backtesting test
(``did the 1-year maximum realised exposure fall below PFE$_{97.5\%}$?'') is
replaced with the following horizon-by-horizon comparison consistent with
the Basel traffic light framework for VaR:

\paragraph{Test Construction.}
At each business day $d$ in the backtesting window, and for each fixed
forward horizon $h$ (1M, 3M, 6M, 1Y), compare the model's predicted
$\text{PFE}_{97.5\%}(h)$ as of day $d$ to the realised netting-set MtM
(net of collateral) at day $d + h$:

\begin{equation}
  \text{Exception at }(d, h):
    \quad \tilde{V}_{\text{realised}}(d+h) > \text{PFE}_{97.5\%}(d, h)
  \label{eq:backtest_exception}
\end{equation}

An exception occurs when the realised collateralised exposure \emph{at that
specific forward horizon} exceeds the predicted PFE. This is conceptually
identical to VaR backtesting (Basel green/yellow/red zone), applied per
horizon rather than globally.

\paragraph{Traffic Light Zones.}

\begin{table}[h!]
\centering
\begin{tabular}{@{}llll@{}}
\toprule
\textbf{Zone} & \textbf{Exceptions (out of 250 obs.\ per horizon)} & \textbf{Expected} & \textbf{Action} \\
\midrule
Green  & 0--6   & $\leq 2.5\%$  & No action required \\
Yellow & 7--9   & $2.5\%$--$3.6\%$ & Explanation required; potential add-on \\
Red    & 10+    & $>3.6\%$      & Model review mandatory; capital add-on \\
\bottomrule
\end{tabular}
\caption{Backtesting traffic light zones per horizon (250 trading days)}
\end{table}

\paragraph{Reporting.}
Backtesting results are reported quarterly to the MRO for each of the four
horizons (1M, 3M, 6M, 1Y) separately. A breach at any horizon triggers the
Yellow or Red zone action at that horizon independently; a Red zone at any
single horizon constitutes a model failure.

\paragraph{Collateral Netting in Realised Exposure.}
Realised exposure $\tilde{V}_{\text{realised}}(d+h)$ is computed as the actual
netting-set MtM at day $d+h$, net of the collateral balance actually held
on that date, consistent with how PFE is modelled. The collateral balance
used is the \emph{actual} balance, not the simulated one, to avoid compounding
simulation error in the backtesting comparison.

\subsection{Stress Scenario Analysis}
\label{sec:stress_scenarios}

Five stress scenarios are applied to the current portfolio quarterly:

\begin{table}[h!]
\centering
\begin{tabular}{@{}lrr@{}}
\toprule
\textbf{Scenario} & \textbf{EE Increase (\$M)} & \textbf{PFE$_{97.5\%}$ Increase (\$M)} \\
\midrule
2008 GFC replay              & +\$1{,}100 & +\$2{,}400 \\
2020 COVID stress            & +\$650     & +\$1{,}500 \\
Rates +200bp parallel shift  & +\$420     & +\$980 \\
Credit spreads +100bp        & +\$280     & +\$620 \\
FX +/-20\% (dollar weakening) & +\$210     & +\$450 \\
\bottomrule
\end{tabular}
\caption{Stress scenario EE and PFE impacts (approximate, current portfolio)}
\end{table}

% ─────────────────────────────────────────────────────────────────────────────
\section{Performance Metrics and Calibration Standards}
% ─────────────────────────────────────────────────────────────────────────────

\begin{longtable}{@{}llll@{}}
\toprule
\textbf{Metric} & \textbf{Acceptable Range} & \textbf{Frequency} & \textbf{Action if Breached} \\
\midrule
\endhead
\bottomrule
\endfoot
EE P\&L Explain                & $>85\%$           & Daily    & Investigate within 2 business days \\
SA-CCR ratio (SA-CCR / 1.4$\times$PFE) & $[0.9,\, 1.5]$  & Quarterly & Review within 30 days \\
Backtest exceptions per horizon & $\leq 6$ (Green)  & Quarterly & Escalate if Yellow; mandatory if Red \\
Stressed EPE multiplier $m_{\text{stress}}$ & $[1.0,\, 3.0]$ & Quarterly & Escalate if $>3.0$ \\
Correlation matrix condition \# & $<1{,}000$        & Annual   & Recalibrate; notify MRO \\
Data completeness (CSA coverage) & $\geq 99\%$       & Daily    & Escalate to Legal/Operations \\
Dispute-history flag review     & N/A               & Weekly   & Flag discrepancy to Collateral Mgmt \\
Intraday run completion time    & $<30$ minutes      & Daily    & Technology escalation \\
Full EOD run completion time    & $<4$ hours         & Daily    & Technology escalation \\
Real-world scaling factor $\xi_{\text{RW}}$ & $[1.0,\, 1.5]$ & Quarterly & MRO notification; re-estimate \\
\end{longtable}

\subsection{Recalibration Schedule}

\begin{itemize}[noitemsep]
  \item \textbf{HW2 and CIR parameters}: Annual recalibration to 5-year
        historical data; ad hoc recalibration upon material market structure
        change (e.g.\ rate regime change).
  \item \textbf{Heston equity parameters}: Quarterly recalibration to current
        vol surface.
  \item \textbf{Correlation matrix}: Annual recalibration using rolling 5-year
        equally-weighted data; reviewed after market stress events.
  \item \textbf{Real-world risk premium $\hat{\phi}(t)$}: Quarterly
        re-estimation using expanded historical window.
  \item \textbf{SA-CCR supervisory parameters}: Fixed by BCBS~279; updated
        only upon regulatory change with full MRO review.
\end{itemize}

% ─────────────────────────────────────────────────────────────────────────────
\section{Limitations, Assumptions, and Known Weaknesses}
\label{sec:limits}
% ─────────────────────────────────────────────────────────────────────────────

\subsection{Gaussian Copula Tail Dependence}

The joint factor simulation uses a Gaussian copula. The Gaussian copula
underestimates tail dependence --- the probability of multiple factors
simultaneously moving to extreme values is lower in the Gaussian copula
than in empirical data, particularly during systemic stress events.

\textbf{Compensating control}: stressed EE (computed using 2008 and 2020 shock
scenarios) is disclosed alongside base EE. A capital add-on for tail dependence
is included in the model reserve.

\subsection{Static Correlation Matrix}

The $59\times 59$ correlation matrix is calibrated once per year. During
periods of market stress, correlations can change rapidly and materially.

\textbf{Compensating control}: the Market Risk Officer monitors intraday
correlation estimates from a 60-day rolling window and can trigger an
unscheduled recalibration if any major pair moves $>30$ percentage points.

\subsection{Equity Gap Risk}

The Heston model uses continuous-time diffusion with no jump component.
Overnight equity gaps are not captured.

\textbf{Compensating control}: equity derivatives with remaining maturity
$<6$ months use a jump-adjusted vol ($\sigma_{\text{adj}} = 1.2 \times
\sigma_{\text{Heston}}$) for PFE calculation only.

\subsection{Proxy Credit Spread Mapping}

Approximately 8\% of counterparties by notional have no observable CDS spread.

\textbf{Compensating control}: proxy-mapped counterparties flagged in all
outputs; 25\% EPE add-on in CVA calculation; quarterly review by Credit Risk.

\subsection{Exotic Derivative Approximations}

Path-dependent exotic derivatives use closed-form approximations or
pre-computed exposure grids, potentially introducing pricing errors of up to 5\%.

\textbf{Compensating control}: exotics exceeding \$50M in notional receive a
full nested Monte Carlo re-price on a monthly basis.

\subsection{FX Volatility Smile (Long-Dated)}
\label{sec:fx_smile}

The GBM FX model uses a single ATM volatility calibrated to the simulation
horizon. For long-dated positions ($>5$ years), this ignores the FX
volatility smile and the tendency for long-dated FX vol surfaces to exhibit
significant wing premium.

\textbf{Quantification}: for EUR/USD FX options with 5--10 year maturity,
the model is estimated to understate PFE$_{97.5\%}$ by 8--15\% relative to
a local-vol or stochastic-vol benchmark. This affects approximately 4\% of
notional in the current portfolio.

\textbf{Compensating control}: a 1.15$\times$ vol scaling add-on is applied
to FX derivatives with maturity $>5$ years for PFE calculation. A SABR or
local-vol FX model is a v2.0 development candidate (Appendix~\ref{app:imm_gap}).

\subsection{Real-World Scaling: Time-Independence}

The real-world scaling factor $\xi_{\text{RW}}$ is applied as a
time-independent scalar across all horizons (see Section~\ref{sec:measure}).
The true measure adjustment is horizon-dependent: at short horizons
($<6$ months) the $\mathbb{Q}$--$\mathbb{P}$ difference is small; at long
horizons ($>5$ years) it can be substantial.

\textbf{Compensating control}: $\xi_{\text{RW}}$ is re-estimated quarterly.
Full $\mathbb{P}$-measure simulation is planned for v2.0.

% ─────────────────────────────────────────────────────────────────────────────
\section{Compensating Controls}
\label{sec:controls}
% ─────────────────────────────────────────────────────────────────────────────

\begin{longtable}{@{}lll@{}}
\toprule
\textbf{Risk} & \textbf{Control} & \textbf{Owner} \\
\midrule
\endhead
\bottomrule
\endfoot
Tail dependence underestimation
  & Stressed EE disclosure; model reserve add-on
  & XVA Quant / MRO \\
Static correlation
  & 60-day rolling monitor; unscheduled recalibration at 30pp move
  & Market Risk \\
Equity gap risk
  & 1.2$\times$ vol adjustment for short-dated equity derivatives
  & XVA Quant \\
Proxy CDS mapping
  & 25\% EPE add-on; quarterly review by Credit Risk
  & Credit Risk \\
Exotic approximation error
  & Monthly full reprice for $>$\$50M; quarterly benchmark
  & XVA Technology \\
FX smile (long-dated)
  & 1.15$\times$ vol scaling for FX derivatives $>5$Y maturity
  & XVA Quant \\
Engine failure (downstream cascade)
  & Fallback: prior-day EPE/ENE scaled by Greeks; alert to desks
  & XVA Technology \\
CSA database gaps
  & Missing CSA defaults to uncollateralised (conservative)
  & Legal / Operations \\
Data staleness
  & 2-business-day staleness alert; fallback to prior-day with flag
  & Market Data \\
Real-world scaling error
  & Quarterly re-estimation of $\xi_{\text{RW}}$; MRO notification if outside $[1.0, 1.5]$
  & XVA Quant \\
Dispute-history MPoR trigger
  & Weekly reconciliation of dispute flags; doubled MPoR applied automatically
  & Collateral Management \\
Short-dated time grid inadequacy
  & Delta-approximation for $<3$M instruments; flagged in credit limit reports
  & XVA Technology \\
\end{longtable}

\subsection{Model Reserve}

A model reserve of \textbf{\$35 million} is held against PFE/CCR model
uncertainty:

\begin{itemize}[noitemsep]
  \item Gaussian copula tail dependence: \$15M;
  \item Proxy CDS mapping basis risk: \$10M;
  \item Exotic derivative path-dependency approximation: \$5M;
  \item FX smile and real-world scaling uncertainty: \$5M.
\end{itemize}

The reserve is reviewed annually alongside the model recalibration cycle and
may be adjusted following material market events.

% ─────────────────────────────────────────────────────────────────────────────
\section{Relationship to Downstream Models}
% ─────────────────────────────────────────────────────────────────────────────

The PFE/CCR Exposure Engine is a producer model; it does not consume outputs
from other models within the XVA suite. The table below summarises the
downstream dependencies:

\begin{table}[h!]
\centering
\begin{tabular}{@{}llll@{}}
\toprule
\textbf{Downstream Model} & \textbf{Model ID} & \textbf{Input Consumed} & \textbf{Sensitivity} \\
\midrule
CVA  & APEX-MDL-0010 & EE profile (all counterparties)        & High --- linear in EE \\
FVA  & APEX-MDL-0011 & EE and ENE profiles                    & High --- linear in EE/ENE \\
MVA  & APEX-MDL-0012 & Simulated IM profiles (Monte Carlo)    & High --- direct path dependency \\
ColVA & APEX-MDL-0013 & $|\text{EE}|$ for collateral balance  & Medium \\
\bottomrule
\end{tabular}
\caption{Downstream model dependencies}
\end{table}

Any change to the PFE/CCR engine that affects EE or ENE profiles triggers
a mandatory impact assessment for all four downstream models before deployment
to production.

% ─────────────────────────────────────────────────────────────────────────────
\section{Change Management}
% ─────────────────────────────────────────────────────────────────────────────

\textbf{Material changes} requiring full MRO re-validation:

\begin{itemize}[noitemsep]
  \item Change to any market factor model (HW2 $\to$ LMM, Heston $\to$ SABR, etc.);
  \item Change to netting set aggregation logic or MPoR treatment;
  \item Change to SA-CCR add-on calculation methodology;
  \item Change to real-world drift adjustment methodology;
  \item Integration of new asset class or new market;
  \item Migration to a new simulation infrastructure or hardware platform.
\end{itemize}

\textbf{Non-material changes} via expedited review (5 business days):

\begin{itemize}[noitemsep]
  \item Annual recalibration of factor model parameters (within pre-approved bounds);
  \item Addition of a new counterparty to the proxy CDS mapping table;
  \item Quarterly re-estimation of real-world scaling factor $\xi_{\text{RW}}$;
  \item Computational optimisation with no change to outputs;
  \item Bug fixes that correct errors without changing methodology.
\end{itemize}

% ─────────────────────────────────────────────────────────────────────────────
\section{Approvals and Sign-Off}
% ─────────────────────────────────────────────────────────────────────────────

\begin{table}[h!]
\centering
\begin{tabular}{@{}llll@{}}
\toprule
\textbf{Role} & \textbf{Name} & \textbf{Status} & \textbf{Date} \\
\midrule
Model Developer (v1.1)   & XVA Quant Team      & $\checkmark$ Submitted & March 2026 \\
Model Risk Officer       & Dr.\ Rebecca Chen   & $\circ$ Pending & --- \\
Chief Risk Officer       & Dr.\ Priya Nair     & $\circ$ Pending & --- \\
Business Owner           & James Okafor        & $\circ$ Pending & --- \\
\bottomrule
\end{tabular}
\caption{Approval status --- v1.1}
\end{table}

\vspace{2em}
\hrule
\vspace{0.5em}
{\small\color{apexgray}
Document Classification: RESTRICTED --- Model Risk Sensitive\\
Apex Global Bank \textbar\ Quantitative Research Division \textbar\ XVA Team\\
Next scheduled review: March 2027 (or upon material model change)
}

% ─────────────────────────────────────────────────────────────────────────────
\appendix
\section{Factor Space Decomposition}
\label{app:factors}
% ─────────────────────────────────────────────────────────────────────────────

\textbf{[NEW]} The following table documents the complete set of stochastic
factors in the joint simulation, correcting the factor count stated as 47
in v1.0. The production implementation has always used 59 factors; the error
was in the documentation.

\begin{longtable}{@{}llll@{}}
\toprule
\textbf{Asset Class} & \textbf{Model} & \textbf{Factors per Instrument} & \textbf{Subtotal} \\
\midrule
\endhead
\bottomrule
\endfoot
IR: SOFR  & HW2 ($x_1, x_2$) & 2 & 2 \\
IR: ESTR  & HW2 ($x_1, x_2$) & 2 & 4 \\
IR: SONIA & HW2 ($x_1, x_2$) & 2 & 6 \\
IR: TONAR & HW2 ($x_1, x_2$) & 2 & 8 \\
IR: HIBOR & HW2 ($x_1, x_2$) & 2 & 10 \\
FX: 10 pairs vs.\ USD & GBM ($\log S$) & 1 each & 20 \\
EQ: 12 equity indices & Heston ($\log S, v$) & 2 each & 44 \\
Credit: 15 spread curves & CIR ($\lambda$) & 1 each & 59 \\
Commodity: 5 indices & GBM ($\log S$) & 1 each & --- \\
\bottomrule
\end{longtable}

\textbf{Note}: commodity indices ($5 \times 1 = 5$ factors) bring the total to 64
if included. In the current production system, commodities are not jointly
simulated via the Cholesky decomposition; their add-on is computed under the
SA-CCR supervisory approach without simulation. The joint Cholesky
decomposition therefore covers $59$ factors ($= 10_{\text{IR}} + 10_{\text{FX}}
+ 24_{\text{EQ}} + 15_{\text{credit}}$). Commodity exposure is captured
via SA-CCR add-on only. A planned v2.0 enhancement will incorporate commodity
simulation into the joint framework.

\subsection*{Cross-Asset Correlation Block Structure}

The $59 \times 59$ correlation matrix $\mathbf{P}$ is structured in blocks:

\begin{itemize}[noitemsep]
  \item \textbf{IR-IR}: $10 \times 10$ block; HW2 within-curve correlation $\rho_{12} = -0.72$;
        cross-curve IR correlations estimated from 5-year historical data.
  \item \textbf{IR-FX}: $10 \times 10$ block; key entries include
        SOFR--EUR/USD ($\rho = -0.31$), SOFR--JPY/USD ($\rho = +0.18$).
  \item \textbf{EQ-EQ}: $24 \times 24$ block; within-index
        $\text{Corr}(\log S, v) = \rho_{\text{SV}}$ per Heston; cross-index
        equity correlations $0.60$--$0.85$ for developed markets.
  \item \textbf{credit-credit}: $15 \times 15$ block; IG-IG within
        class $\sim 0.70$; HY-HY $\sim 0.55$.
  \item \textbf{EQ-credit}: $24 \times 15$ block; equity-to-single-name
        spread correlations $\sim -0.35$ (risk-off relationship).
\end{itemize}

% ─────────────────────────────────────────────────────────────────────────────
\section{IMM Approval Gap Analysis}
\label{app:imm_gap}
% ─────────────────────────────────────────────────────────────────────────────

\textbf{[NEW]} This appendix provides the gap analysis between the current
model and Basel IMM requirements (BCBS SA-CCR rules, as codified in
12 CFR Part 3 for US banks), required as a prerequisite for IMM use approval.

\begin{longtable}{@{}lllll@{}}
\toprule
\textbf{IMM Requirement} & \textbf{Ref.} & \textbf{Current Status} & \textbf{Gap} & \textbf{Remediation} \\
\midrule
\endhead
\bottomrule
\endfoot
Full real-world measure for PFE
  & BCBS \S\,289
  & Partial (scaling factor)
  & Material
  & Full $\mathbb{P}$ sim, v2.0 \\
Stressed EPE using stressed params
  & BCBS \S\,277
  & Implemented (v1.1)
  & Closed
  & --- \\
Daily margining with MPoR $\geq 10$ days
  & BCBS 261
  & Implemented
  & Closed
  & --- \\
Dispute-history MPoR doubling
  & BCBS 261 \S49
  & Implemented (v1.1)
  & Closed
  & --- \\
Netting set-level EPE
  & BCBS \S\,284
  & Implemented
  & Closed
  & --- \\
Backtesting vs.\ realised exposure
  & BCBS \S\,302
  & Revised (v1.1)
  & Substantially closed
  & Annual audit \\
Effective EPE for capital
  & BCBS \S\,279
  & Implemented
  & Closed
  & --- \\
Quarterly stressed calibration
  & BCBS \S\,278
  & Implemented (v1.1)
  & Closed
  & --- \\
Joint simulation of IR, FX, EQ, credit
  & BCBS \S\,290
  & Implemented (59 factors)
  & Closed
  & --- \\
Commodity simulation (not SA-CCR)
  & BCBS \S\,290
  & Not implemented
  & Open
  & v2.0 scope \\
FX smile model for long-dated
  & BCBS \S\,293
  & Partial (vol add-on)
  & Minor
  & v2.0 SABR/LV \\
Daily time grid (vs.\ weekly)
  & BCBS \S\,284
  & Weekly to 1Y, then monthly
  & Minor ($<3$M instruments)
  & v2.0 \\
Independent model validation
  & SR 11-7 \S\,V
  & In progress (Dr.\ Chen)
  & Procedural
  & H2 2026 \\
1-year parallel run requirement
  & 12 CFR 3.132(d)(3)
  & Not started
  & Critical (pre-IMM)
  & Commence Q3 2026 \\
\end{longtable}

\textbf{IMM readiness assessment}: the model is approximately 70\% complete
relative to full IMM requirements. The two critical open items are the
full real-world measure implementation and the 1-year parallel run requirement
(which must precede any IMM use approval application). Estimated timeline to
IMM readiness: 18--24 months from March 2026, contingent on v2.0 delivery
and regulatory review timeline.

% ─────────────────────────────────────────────────────────────────────────────
\section{Stressed EPE Calibration Details}
\label{app:stressed_epe}
% ─────────────────────────────────────────────────────────────────────────────

\textbf{[NEW]} This appendix documents the identification of the stressed
calibration period and the resulting parameter shifts.

\subsection*{Stressed Period Identification}

The stressed period must cover a ``significantly stressed'' 12-month window
per BCBS~\S\,277. The following candidates were evaluated:

\begin{table}[h!]
\centering
\begin{tabular}{@{}lrr@{}}
\toprule
\textbf{Candidate Period} & \textbf{EPE Stress Multiplier} & \textbf{Coverage of Portfolio Risk Factors} \\
\midrule
Oct 2008 -- Sep 2009 (GFC)   & 1.42$\times$ & High (IR, credit, FX, EQ) \\
Mar 2020 -- Feb 2021 (COVID) & 1.28$\times$ & Medium (EQ, credit; IR less stressed) \\
Jun 2022 -- May 2023 (rate hikes) & 1.19$\times$ & Medium (IR; EQ moderate) \\
\bottomrule
\end{tabular}
\caption{Stressed period candidates and multipliers}
\end{table}

\textbf{Selected period}: Oct 2008 -- Sep 2009, as it produces the highest
EPE stress multiplier and covers all major risk factors. This selection is
documented and approved by the Market Risk Officer annually.

\subsection*{Stressed Parameter Shifts}

\begin{table}[h!]
\centering
\begin{tabular}{@{}llll@{}}
\toprule
\textbf{Parameter} & \textbf{Base} & \textbf{Stressed} & \textbf{Impact} \\
\midrule
$\sigma_1^{\text{HW2}}$ (IR vol, factor 1) & 0.007 & 0.018 & $+2.6\times$ \\
$\sigma_2^{\text{HW2}}$ (IR vol, factor 2) & 0.014 & 0.031 & $+2.2\times$ \\
$\sigma^{\text{FX}}_{\text{EUR/USD}}$ & 0.075 & 0.142 & $+1.9\times$ \\
$\xi^{\text{Heston}}$ (EQ vol-of-vol) & 0.35  & 0.71  & $+2.0\times$ \\
$\sigma^{\text{CIR}}_\lambda$ (credit vol) & 0.45 & 1.20 & $+2.7\times$ \\
\bottomrule
\end{tabular}
\caption{Selected parameter shifts: base vs.\ GFC stressed calibration}
\end{table}

% ─────────────────────────────────────────────────────────────────────────────
\section{Notation Summary}
\label{app:notation}
% ─────────────────────────────────────────────────────────────────────────────

\begin{longtable}{@{}ll@{}}
\toprule
\textbf{Symbol} & \textbf{Definition} \\
\midrule
\endhead
\bottomrule
\endfoot
$V_k(t)$          & MtM value of trade $k$ at time $t$ (Apex's perspective) \\
$V_{\text{net}}(t)$ & Net MtM value of netting set at time $t$ \\
$C(t)$            & Collateral balance at time $t$ (positive = held by Apex) \\
$C_{\text{net}}(t)$ & Haircut-adjusted collateral balance \\
$C^*(t, \omega)$  & Target collateral (CSA credit support amount) on path $\omega$ \\
$C_{\text{frozen}}(t, \omega)$ & MPoR-frozen collateral balance \\
$\Theta$          & CSA threshold \\
$\delta$          & CSA minimum transfer amount (MTA) \\
$\tau$            & Margin Period of Risk (in time units) \\
$\Delta_s$        & Settlement lag (business days) \\
$\tilde{V}(t)$    & MPoR-adjusted collateralised exposure \\
$\text{EE}(t)$    & Expected Exposure at time $t$ (risk-neutral) \\
$\text{EPE}$      & Time-averaged Expected Positive Exposure (1-year) \\
$\text{EffEE}(t)$ & Effective Expected Exposure (non-decreasing) \\
$\text{EffEPE}$   & Effective EPE (Basel IMM capital metric) \\
$\text{ENE}(t)$   & Expected Negative Exposure at time $t$ (risk-neutral) \\
$\text{PFE}_\alpha(t)$ & Potential Future Exposure at confidence $\alpha$ (real-world) \\
$\xi_{\text{RW}}$ & Real-world PFE scaling factor ($= 1.12$ as of March 2026) \\
$\phi(t)$         & Real-world drift adjustment function \\
$\alpha$          & SA-CCR supervisory scalar ($= 1.4$) \\
$\text{RC}$       & SA-CCR Replacement Cost \\
$\text{EAD}$      & Exposure at Default (SA-CCR) \\
$d_j$             & SA-CCR adjusted notional (supervisory duration for IR) \\
$\text{SD}_j$     & Supervisory duration for trade $j$ \\
$\text{MF}_j$     & SA-CCR maturity factor for trade $j$ \\
$D_b^{\text{ccy}}$ & Effective bucket notional for IR bucket $b$, currency ccy \\
$\rho_{bb'}$      & SA-CCR IR bucket correlation ($\rho_{12}=\rho_{23}=0.70$; $\rho_{13}=0.30$) \\
$N$               & Number of Monte Carlo paths ($= 10{,}000$) \\
$M$               & Number of simulation time steps ($= 200$) \\
$a_1, a_2$        & Hull-White mean-reversion speeds \\
$\sigma_1, \sigma_2$ & Hull-White volatilities \\
$\kappa, \bar{v}, \xi$ & Heston mean-reversion, long-run variance, vol-of-vol \\
$\kappa_\lambda, \bar{\lambda}, \sigma_\lambda$ & CIR credit spread parameters \\
$m_{\text{stress}}$ & Stressed EPE multiplier ($= 1.42$, GFC calibration) \\
$\mathbf{P}$      & $59 \times 59$ factor correlation matrix \\
$\mathbf{L}$      & Cholesky factor of $\mathbf{P}$ \\
$\mathbb{Q}$      & Risk-neutral probability measure \\
$\mathbb{P}$      & Real-world (physical) probability measure \\
\end{longtable}

% ─────────────────────────────────────────────────────────────────────────────
\end{document}
